\chapter{Introducción}
\label{ch:introduccion}

% Contexto general del proyecto

El crecimiento exponencial de los deportes electrónicos (esports) en la última década ha generado un volumen masivo de información distribuida en múltiples plataformas digitales. En particular, el ecosistema competitivo de \textit{Counter-Strike}, uno de los juegos más longevos y populares en la escena profesional, cuenta con cientos de equipos, miles de jugadores y un constante flujo de fichajes, resultados de torneos y cambios en las plantillas. Esta información, aunque valiosa para aficionados, analistas y profesionales del sector, se encuentra fragmentada en diversas fuentes como sitios web especializados, redes sociales y wikis colaborativas.

% Motivación del trabajo

Liquipedia se ha consolidado como la enciclopedia colaborativa de referencia en esports, mantenida por la comunidad y estructurada de forma similar a Wikipedia. Sin embargo, el acceso a esta información requiere navegación manual entre múltiples páginas, lo que dificulta la obtención rápida de respuestas a consultas específicas. En este contexto, los sistemas de pregunta-respuesta conversacionales basados en Procesamiento del Lenguaje Natural (PLN) representan una oportunidad para facilitar el acceso a información estructurada mediante interfaces intuitivas en lenguaje natural.

El presente trabajo aborda el desarrollo de un sistema completo de minería de texto y PLN que integra técnicas de web scraping, preprocesamiento lingüístico, traducción automática y generación aumentada por recuperación (RAG, \textit{Retrieval-Augmented Generation}). El sistema resultante permite a los usuarios consultar información sobre equipos profesionales de \textit{Counter-Strike} en español de forma conversacional, superando las barreras idiomáticas y de navegación que presentan las fuentes originales.

% Objetivos principales

Los objetivos principales de este proyecto son:

\begin{itemize}
    \item Construir un corpus en español sobre equipos profesionales de \textit{Counter-Strike} mediante web scraping automatizado de Liquipedia.
    \item Implementar un pipeline de preprocesamiento que incluya limpieza, normalización y traducción automática de textos.
    \item Desarrollar un sistema RAG que combine búsqueda semántica con modelos de lenguaje para generar respuestas contextualizadas.
    \item Analizar los sesgos presentes en el corpus y proponer estrategias de mitigación.
    \item Evaluar el rendimiento del sistema mediante métricas estándar de recuperación y generación.
\end{itemize}

% Estructura de la memoria

% El resto de la memoria se organiza de la siguiente manera: el Capítulo~\ref{ch:estado_arte} revisa los fundamentos teóricos y el estado del arte en sistemas RAG, web scraping y traducción automática; el Capítulo~\ref{ch:corpus_datos} describe el proceso de construcción del corpus mediante scraping de Liquipedia; el Capítulo~\ref{ch:eda} presenta un análisis exploratorio de los datos recopilados; el Capítulo~\ref{ch:sesgos} examina los sesgos identificados en el corpus y las medidas adoptadas; el Capítulo~\ref{ch:metodologia} detalla la arquitectura del sistema y las técnicas implementadas; el Capítulo~\ref{ch:resultados} expone los resultados de la evaluación; y finalmente, el Capítulo~\ref{ch:discusion} discute las limitaciones del trabajo y las líneas futuras de investigación.

